% Round One (NeurIPS poster) submission, Smart Buildings Challenge 2026.
%
% Written against plain article rather than the NeurIPS workshop style file,
% which is not public at the time of writing. Swap \documentclass and the
% geometry line for the official template before submitting; nothing else in
% this file depends on the class.
\documentclass[10pt]{article}

\usepackage[margin=1in]{geometry}
\usepackage{amsmath}
\usepackage{amssymb}
\usepackage{booktabs}
\usepackage{microtype}
\usepackage{url}
\usepackage[colorlinks=true,allcolors=blue]{hyperref}
\usepackage{caption}

\title{\textbf{Airflow Is the Load, Not Just the Bill}\\[2pt]
\large A Mispriced Actuator in the Smart Buildings Challenge}

\author{%
  Minh Nhat Nguyen\\
  \texttt{github.com/MinhNhat-2504/hvac-generalization-lab}
}

\date{}

\begin{document}
\maketitle

\begin{abstract}
The Smart Buildings Challenge asks whether a single HVAC control policy can
transfer to buildings and climates it has never seen, and poses as an open
question whether domain-specific models or generic foundation models will
generalize best. We report a result on the domain-specific side. All four
reference controllers shipped with the challenge price fan static pressure as
a cost to be paid in proportion to demand. In this plant that is the wrong
sign: the simulator caps recirculation at 70\%, so three tenths of every cubic
metre the fan moves is ambient air that must then be conditioned from the
outdoor temperature. Raising static pressure to deliver heat therefore
\emph{increases} the heating load it was raised to meet. Correcting this single
term---clamping static pressure to its floor whenever an air handler is
heating---reduces the leaderboard metric from $0.05496$ to $0.03682$ over the
30 published training scenarios, a $33.0\%$ reduction, and improves comfort,
energy cost and carbon \emph{simultaneously} rather than trading between them.
On the worst scenario the metric falls $66\%$. The controller has no learned
parameters, runs in pure NumPy, and is a 40-line change to the strongest
shipped baseline. We also report five refuted hypotheses, including our own
initial one, and a methodological finding: the baseline return table shipped
with the starter kit predates a simulator fix and no longer reproduces.
\end{abstract}

\section{Introduction}

The challenge \cite{goldfeder2026} evaluates one control policy on six unseen
buildings at five unseen weather stations. Buildings differ in floor plan, zone
count, air-handler count and envelope construction, so both the observation and
action vectors change shape per building and there is no fixed layout to
memorize. The organizers state the core hurdle as transferability, and ask
whether ``domain specific models will win out, or whether generic foundation
models will generalize best'' \cite{goldfeder2026}.

Two features of the setup push against learned policies. Per-rollout wall-clock
is bounded at ten times the \texttt{bang\_bang} baseline, which rules out
winning by online planning. And the metric is a mean over 30 equally weighted
scenarios whose difficulty is very unequal: in our measurements the four worst
scenarios contribute $38\%$ of the strongest baseline's total loss. Mean-case
improvement is worth a fraction of tail-case improvement, and tails are where
sample-efficient learning is weakest.

This paper takes the other route. We read the reward function and the plant
model as a price list, find one actuator whose price the reference controllers
have the wrong sign for, and correct it.

\section{The reward as a price list}

Per five-minute step the reward is
\begin{equation}
r_t = 0.6\,p_t - 0.2\,c_t - 0.2\,e_t ,
\end{equation}
with all three terms penalties, $r_t \le 0$, and the leaderboard metric the mean
over scenarios of $|\bar r|$. Substituting the simulator's own cost and carbon
models at the shipped normalisation constants
($320$\,kW electricity, $820$\,kW gas) gives the marginal cost in metric units
of one delivered kilowatt-hour:

\begin{center}
\begin{tabular}{lc}
\toprule
 & metric cost per kWh \\
\midrule
fan / pump electricity            & $0.0064$ \\
gas heat (boiler, $\eta\approx1$) & $0.0034$ \\
chiller cooling (COP $3.5$)       & $0.0018$ per kWh \emph{thermal} \\
\bottomrule
\end{tabular}
\end{center}

Electricity costs $1.9\times$ gas per kWh of fuel, and a fan has no coefficient
of performance to divide by where a chiller has $3.5$. \textbf{Moving air is the
most expensive thing this plant can do.}

The comfort term is an occupancy-weighted mean of per-zone shortfall, where
shortfall is $\mathrm{dev}/5\,\mathrm{K}$ inside the knee and quadratic and
unbounded beyond it. Every occupied zone sitting $1$\,K outside its band
therefore costs $0.6 \times 0.2 = 0.12$ per step, which at $0.0064$ per kWh is
$19$\,kWh per step---$225$\,kW held continuously, against a $320$\,kW plant.
\emph{While comfort is scored, energy is nearly free beside it.}

Finally, the comfort term is gated hard on occupancy:
\verb|total_occupancy > 0.0| or the term is identically zero whatever the zone
temperatures are. Occupants arrive between hours 6 and 12, leave between 13 and
19, and are absent on weekends and holidays, leaving roughly $60\%$ of the
$8{,}640$ scored steps with the comfort term switched off.

\section{Airflow is the load}

The gated comfort term invites an obvious move: stop conditioning an empty
building. We implemented it and it was \emph{worse} than the strongest baseline
on every scenario we tried. Section~\ref{sec:negative} gives the two reasons.
The useful finding came from asking instead why the reference controllers spend
what they spend.

Every policy in the shipped \texttt{pid} family sets static pressure from the
control effort $u$:
\begin{equation}
\texttt{pressure} = \mathrm{clip}(2|u| - 1,\,-1,\,1).
\label{eq:pid}
\end{equation}
The baseline's own documentation defends this against pinning pressure at
maximum, and against that alternative it is an improvement. But $|u|$ does not
distinguish the two directions of the loop, and the directions are not
symmetric.

\paragraph{Heating.} The air handler cannot heat. It mixes return and outside
air, cools the mixture towards the commanded supply temperature, and all heat
in the building comes from the VAV reheat coils fed by the hot-water loop. In
cold weather the supply-air action is inert as well: the handler delivers
$\min(\text{mixed air}, \text{setpoint})$, and with recirculation capped at
$0.70$, return air near $293$\,K and ambient at $231$\,K,
\begin{equation}
\text{mixed} = 0.70 \times 293 + 0.30 \times 231 = 274\ \mathrm{K},
\end{equation}
already below the $285$\,K floor the setpoint can command.

So in heating the fan moves $274$\,K air past a coil, and every cubic metre of
it must be heated from $274$\,K before it does the zone any good. The
ventilation floor makes this unavoidable: three tenths of the moved air is
ambient, at $-42^\circ$C in the coldest shipped strip. \textbf{The fan work is
spent, and then gas is spent undoing the air the fan moved.}

The thermal energy a zone needs is fixed by its heat loss. It can arrive as a
lot of slightly warm air or a little hot air, and the second is strictly cheaper
here, because the gas is the same either way and the fan is not. Equation
\eqref{eq:pid} does the opposite: it buys more of the expensive fuel in order to
import more of the load.

\paragraph{Cooling.} Air is the only delivery mechanism, so flow must be bought.
But the trade still favours temperature over flow---the chiller kilowatt-hour is
a third of the fan kilowatt-hour---and the same ventilation floor means that in
$45^\circ$C ambient, extra flow again imports extra load.

\section{Controller}

Our policy is the shipped \texttt{pid\_water} controller with two changes and no
learned parameters.

\paragraph{A flat band.} \texttt{pid\_water} widens its band to an eco setback
outside comfort hours. We hold the comfort band around the clock. This is not
an energy argument but a tail argument: a building held in band overnight starts
the morning with no deviation to recover, and the recovery penalty is larger
than the holding cost. The shipped \texttt{fixed\_setpoint} already hints at
this by beating \texttt{pid\_water} in cold climates.

\paragraph{A direction-aware pressure law.} Where \eqref{eq:pid} applies, we
instead clamp static pressure to its floor whenever the handler is heating
($u>0$), and scale it by $0.6$ when cooling. The simulator enforces a minimum
static pressure under any running handler, so the floor requests that minimum
rather than no air; ventilation is not something a policy may trade away, and
this does not try to.

One further change follows from the same reading. \texttt{pid\_water} opens the
reheat valve in proportion to shortfall, fully open at $2$\,K. That protects a
mild climate from overshoot and starves a cold one, where the valve is the only
heat source and the air reaching it is already $20$\,K below the zone. We
saturate the valve at $0.5$\,K.

\section{Results}

All runs are full $8{,}640$-step episodes. The simulator is deterministic---we
verified two runs of one policy on one scenario reproduce to ten decimal
places---so each number is a single rollout and no averaging is required.

\begin{table}[h]
\centering
\caption{Leaderboard metric (lower is better) over the 30 published training
scenarios, with the weighted contribution of each reward term. Both policies
re-run with the current simulator.}
\begin{tabular}{lccccc}
\toprule
policy & metric & worst scenario & comfort & energy & carbon \\
\midrule
\texttt{pid\_water} (strongest baseline) & $0.05496$ & $0.21392$ & $0.03489$ & $0.01150$ & $0.00857$ \\
\textbf{ours} & $\mathbf{0.03682}$ & $\mathbf{0.07245}$ & $0.02694$ & $0.00572$ & $0.00416$ \\
\midrule
change & $-33.0\%$ & $-66.1\%$ & $-22.8\%$ & $-50.3\%$ & $-51.5\%$ \\
\bottomrule
\end{tabular}
\label{tab:main}
\end{table}

\textbf{All three terms improve at once.} That is the signature of a mispriced
actuator rather than a trade: less static pressure is less airflow, which is
less ambient air imported, which is less load, so the coils keep up where
before they could not. Comfort improves \emph{because} the energy was cut, not
despite it.

The gain is concentrated in the tail, which is what the metric rewards. Our
policy is better on 19 of 30 scenarios; summed over scenarios it improves
$0.575$ and regresses $0.031$, a ratio of $19{:}1$, and the median regression
among the eleven is $0.001$. The four largest improvements are
\texttt{fairbanks\_1} $0.214\!\to\!0.072$, \texttt{phoenix\_4}
$0.160\!\to\!0.043$, \texttt{phoenix\_5} $0.136\!\to\!0.049$ and
\texttt{singapore\_4} $0.110\!\to\!0.024$: the two coldest and the two hottest,
which is what a mechanism acting through the ventilation floor predicts.

\begin{table}[h]
\centering
\caption{Ablation, mean over \texttt{fairbanks\_1}, \texttt{singapore\_6} and
\texttt{phoenix\_4}. The flat band and the pressure law are separable, and the
pressure law carries the result.}
\begin{tabular}{lcl}
\toprule
 & metric & \\
\midrule
\texttt{pid\_water}                     & $0.13478$ & shipped baseline \\
\quad + flat band                       & $0.12321$ & band alone \\
\quad + cooling clamp only              & $0.11315$ & \\
\quad + heating clamp only              & $0.05226$ & \textbf{the single biggest term} \\
\quad + both clamps                     & $0.04967$ & \\
\quad + reheat saturation at $0.5$\,K   & $\mathbf{0.04730}$ & final \\
\bottomrule
\end{tabular}
\label{tab:ablation}
\end{table}

\paragraph{Held-out buildings and climates.} Because the published 30 scenarios
are a training set, we built two replacements from material they never use. The
generator is procedural and ten of the 48 shipped weather strips are unused by
the published set, with a further station, \texttt{minneapolis}, used by none of
them. \textbf{dev} is four newly generated buildings crossed with the five
unused strips at seen stations; \textbf{test} is four further buildings crossed
with \texttt{minneapolis}, so neither the building nor the climate has been
seen. All eight buildings fall inside the $28$--$85$ zone and $3$--$9$
air-handler range the organizers state for the evaluation set. Results on these
sets are reported in the camera-ready; the method has no fitted parameters to
carry over, which is the reason to expect them to hold.

\section{What did not work}
\label{sec:negative}

We report these because each was a reasonable inference from the reward and each
was wrong, which is itself information about the benchmark.

\paragraph{Switching the plant off when the building is empty.} Our original
hypothesis, and the sign was backwards. Two independent reasons. First, being
empty \emph{inside} the scored window is not an opportunity: arrivals are
independent stochastic draws, so the building is routinely empty for a few steps
after hour 6, and because regret is normalized by \emph{total} occupancy the
first person through the door carries their zone's entire comfort score alone.
Second, a controller that coasts must later recover, and the recovery lands in
the low-occupancy morning where each occupant's weight is largest. Driving to
the comfort band overnight to avoid that cost \emph{more} than coasting saved:
unoccupied-step loss rose from $447$ to $830$ for a comfort gain of $0.010$.

\paragraph{Deeper setbacks.} Monotonically worse: $0$\,K $0.123$, $3$\,K
$0.166$, $5$\,K $0.160$, $8$\,K $0.183$. Steady holding beats cycling here.

\paragraph{Aiming inside the comfort band.} The penalty is flat zero anywhere
inside the band, so a margin looked like free insurance against drift. It is
clearly harmful ($0.236$ against $0.123$), and harmful on the \emph{comfort}
term, not merely the energy term. Narrowing a $3$\,K band to $2$\,K leaves the
loop, which saturates at $1/k_p = 2$\,K, oscillating across the real edge.

\paragraph{Forcing the water loop hotter.} Having removed the airflow, raising
the water temperature to compensate looked necessary. It is not: carbon rose
from $0.0134$ to $0.0162$. With the flow down, the baseline's own demand reset
is already hot enough.

\paragraph{Suspending the fan economy when comfort is at stake.} Against the
$35{:}1$ price ratio this should pay---a quiet fan delivers heat slowly, and
slow is only cheap when nobody is being charged. It does not pay ($0.074$
against $0.050$), and the reason is the mechanism itself: the fan is what
\emph{causes} the load, so spending the expensive fuel to deliver heat imports
more cold air than the extra flow delivers.

\paragraph{Capping the hot-water pump.} Also electricity, also pinned at
maximum by the baseline, and also worse when capped ($0.051$ against $0.050$).

\section{A note on the shipped baseline table}

The starter kit ships \texttt{baseline\_returns.json}, and
\texttt{score\_leaderboard.py} prints it beside a participant's score so the
four baselines need not be re-run. We find it no longer reproduces:
\texttt{pid\_water} on \texttt{building\_train\_1/fairbanks\_1} gives
$0.21392$ with current code against the recorded $0.21168$, and $0.05496$
against $0.05414$ over the 30 scenarios.

The file was committed with the initial starter kit and never regenerated,
while a later commit fixed the simulation tracker querying occupancy with VAV
ids---a fix whose own message notes that it ``creates fictitious populations
and consumes random draws, changing subsequent simulation trajectories''. The
trajectories moved and the table did not. The discrepancy is $1.1$--$1.5\%$,
which is small but is the same order as several of the tuning effects above.
Participants comparing against the shipped table are measuring against a
simulator that no longer exists.

\section{Discussion}

On the organizers' open question, this is evidence that domain structure is
worth a great deal here, and cheaply. The result is a 40-line change,
dimensionless, with no fitted constants and nothing to transfer, and it holds
across both temperature extremes because it acts through a constraint---the
$70\%$ recirculation cap---that is a property of the plant rather than of any
building.

We would not read it as evidence against learned policies in general. We would
read it as a statement about what a learned policy would have to discover: not a
smooth function of the observation, but the sign of a coupling between an
actuator and the load it is meant to serve, visible only in the regime where the
supply-air action has gone inert. A policy-gradient method optimizing an
$8{,}640$-step return at ten times \texttt{bang\_bang}'s runtime budget has a
hard road to that. The natural synthesis is to let structure carry the
mechanism and let learning carry what structure cannot express---which is the
direction we intend for the final submission.

\paragraph{Limitations.} The tuning in Table~\ref{tab:ablation} beyond the
heating clamp is small, measured on three scenarios, and we do not claim those
constants transfer. The cooling clamp in particular is a climate trade: at a
gain of $0.3$ it helps \texttt{phoenix} and costs \texttt{singapore}. We have
not tested on the organizers' held-out set, and the public leaderboard scores
only a rotating subset.

\paragraph{Reproducibility.} Code, every sweep in this paper as JSON, and the
held-out generator are at
\url{https://github.com/MinhNhat-2504/hvac-generalization-lab}. Each negative
result above is a registered policy variant, runnable by name.

\begin{thebibliography}{9}
\bibitem{goldfeder2026}
J.~Goldfeder, A.~Rai, P.~M.~Wyder, N.~Li, V.~Dean, D.~Vakalis, J.~Lin, B.~Dong,
M.~Rossetti, N.~Kutz, S.~Arisaka, J.~Sipple, H.~Lipson, and Y.~Bengio.
\newblock The Smart Buildings Challenge: Learning Foundation Models for
Real-World Optimization at Scale.
\newblock arXiv:2610.18511, 2026.

\bibitem{starterkit}
Smart Buildings 2026 Organizers.
\newblock Smart Buildings Competition Starter Kit.
\newblock \url{https://github.com/smart-buildings-2026/starter-kit}, 2026.
\end{thebibliography}

\end{document}
